\begin{pastexam}[label={px:gk3-098-ds-03}]{高考三級資訊處理 98 年 第三題}
（申論；24 分）

假設有下列數種排序方法：(A)bubble sort\quad (B)quick sort\quad (C)heap sort\quad (D)merge sort\quad
(E)radix sort\quad (F)insertion sort。回答下列問題時，請分別以 ABCDEF 之代號答之。（每小題 6 分共 24 分）
\begin{itemize}[leftmargin=2.2em]
\item[㈠] 一個排序法，在輸入資料中有多筆相同資料時，於排序前與排序後，任兩筆相同的資料前後順序不變者，
稱之為「穩定排序法」。請問那些排序法為穩定排序法？
\item[㈡] 假設輸入資料有 n 個，在最糟情形下，那些排序法的時間複雜度為 O（nlogn）？
\item[㈢] 排序程式實作時，那些排序法需要額外的陣列或鏈結串列？
\item[㈣] 在程式實作時，一般使用陣列進行排序。有些時候也需要對鏈結串列進行排序。那些排序法無法對單向
鏈結串列（linearly linked list）進行排序？
\end{itemize}
\end{pastexam}

\begin{pastsoln}
答案：㈠ A、D、E、F；㈡ C、D；㈢ D、E；㈣ C。

先把六種排序法的性質列成一張表（資結書第~14.7~節的比較表；基數排序見資結書第~15.2~節），四個小題都從
這張表讀出來。 $n$ 是筆數，基數排序的鍵有 $d$ 位、每位 $r$ 種值。
\begin{center}
\small
\begin{tabular}{lcccl}
\toprule
排序法 & 最壞時間 & 穩定 & 額外空間 & 要不要隨機存取 \\
\midrule
(A) 氣泡 & $O(n^2)$ & 是 & $O(1)$ & 不用 \\
(B) 快速 & $O(n^2)$ & 否 & 遞迴堆疊 $O(\log n)$--$O(n)$ & 不用（見㈣） \\
(C) 堆積 & $O(n \log n)$ & 否 & $O(1)$ & 要，靠 $i$ 算出 $2i$、$2i + 1$ \\
(D) 合併 & $O(n \log n)$ & 是 & $O(n)$ 的暫存陣列 & 不用 \\
(E) 基數 & $O(d(n + r))$ & 是 & $r$ 個桶子或輸出陣列 & 不用 \\
(F) 插入 & $O(n^2)$ & 是 & $O(1)$ & 不用 \\
\bottomrule
\end{tabular}
\end{center}

\textbf{㈠ 穩定：A、D、E、F}。
\begin{itemize}
\item (A) 氣泡排序只在 $A[j] > A[j+1]$ 時交換相鄰兩筆，相等的兩筆不會交換，相對順序不變。
\item (F) 插入排序把新的一筆往前插時，遇到 $\le$ 它的就停，不會越過相等的資料。
\item (D) 合併時左右兩段的首筆相等，先取左段的（$\le$ 時取左），左段原本就在前面。
\item (E) 基數排序（LSD）每一位都用穩定的分配（依序放進桶子、依序收回），所以整體穩定；它也必須
穩定，否則低位排好的順序會被打亂。
\item (B) 不穩定：分割時遠距離交換。 例如 Lomuto 分割 $\langle 2_a, 2_b, 1\rangle$，基準是 $1$，
$1$ 和 $2_a$ 交換得 $\langle 1, 2_b, 2_a\rangle$，兩個 $2$ 的順序反了。
\item (C) 不穩定：把堆積頂端換到最後時跳過中間的資料。 例如 $\langle 1, 2_a, 2_b\rangle$ 建最大堆積
得 $\langle 2_a, 1, 2_b\rangle$（兩個子節點一樣大，換左邊的 $2_a$），取出 $2_a$ 放到最後，結果
$2_b$ 在 $2_a$ 前面。
\end{itemize}

\textbf{㈡ 最壞 $O(n \log n)$：C、D}。 堆積排序建堆 $O(n)$，之後 $n - 1$ 次下沉各 $O(\log n)$；合併
排序 $T(n) = 2T(n/2) + \Theta(n) = \Theta(n \log n)$，和輸入無關。 (A)(F) 遇到反序輸入是
$\Theta(n^2)$；(B) 每次基準都選到最大或最小值時是 $\Theta(n^2)$。 (E) 的時間 $O(d(n + r))$ 和 $n \log n$
沒有一定的大小關係（$d$ 由鍵的範圍決定，鍵值到 $n^c$ 時 $d = \Theta(\log n)$，才是 $\Theta(n \log n)$），
題目問的是比較排序那一類的保證，不選 E。

\textbf{㈢ 需要額外的陣列或鏈結串列：D、E}。 合併排序合併兩段時要一個大小 $n$ 的暫存陣列；基數排序要
$r$ 個桶子（通常用佇列或鏈結串列做），或計數排序用的輸出陣列。 A、C、F 都是就地排序，只用 $O(1)$ 個
變數；B 也是就地分割，只用遞迴的堆疊，不需要額外的陣列。

\textbf{㈣ 無法對單向鏈結串列排序：C}。 堆積排序把陣列當成完整二元樹，靠下標 $i$ 直接算出子節點 $2i$、
$2i + 1$ 和父節點 $\lfloor i/2 \rfloor$，每次下沉都要跳著存取。 單向鏈結串列只能從頭往後走，走到第 $i$
個就要 $O(i)$，沒辦法照原本的方式做（硬做的話時間不再是 $O(n \log n)$）。 其他五種都能在單向鏈結串列
上做：
\begin{itemize}
\item (A) 只比較、交換相鄰的兩個節點，從頭往後走就行。
\item (F) 每次從原串列取下一個節點，從新串列的頭往後找插入位置（不必往回走），仍是 $O(n^2)$。
\item (D) 合併排序最適合鏈結串列：用快慢指標找中點切開，合併時只改指標，不需要暫存陣列。
\item (E) 基數排序把節點依某一位接到 $r$ 個桶子串列的尾端，再依序串回來。
\item (B) Lomuto 分割只需要兩個都往後走的指標；或者把每個節點依和基準比較的結果接到「小於」「大於等於」
兩條串列，再遞迴、串接。
\end{itemize}
有些參考書把 B 也算進來，理由是 Hoare 分割要從右邊往左掃、單向串列做不到。 那是某一種寫法的限制，
不是快速排序本身的限制，所以本書只答 C。
\end{pastsoln}
